% 由 工程文件\tools\gen_tables.py 生成，勿手改。
\section{术语与记号对照表}

下表给出正文中反复出现的术语。中文译名在全书统一；英文术语在正文首次出现时括注，其后从略。

\subsection*{表 A1　术语}
\begin{longtable}{p{0.22\textwidth}p{0.27\textwidth}p{0.41\textwidth}}
\toprule 中文 & 英文 & 说明 \\ \midrule \endfirsthead
\toprule 中文 & 英文 & 说明 \\ \midrule \endhead
\bottomrule \endfoot
A-模型 / B-模型 & A-model / B-model & 二维拓扑场论的两族：前者依赖辛结构（计数伪全纯曲线），后者依赖复结构（复结构变形的积分）。 \\
镜像对称 & mirror symmetry & 一对流形上 A-模型与 B-模型之间的等价；不是经典几何的等同。 \\
拓扑扭转 & topological twist & 把超对称理论的 Lorentz 群换成与 R-对称的对角子群，使某个超荷成为标量，从而得到拓扑场论。 \\
复辛流形 & complex symplectic manifold & 带全纯非退化闭 2-形式的复流形；典范线丛自动平凡。 \\
超凯勒流形 & hyper-Kähler manifold & 带三个满足四元数关系的复结构、且三个 Kähler 形式均闭的流形。 \\
平坦丛 & flat bundle & 带可积联络的丛；等价于基本群的表示（模共轭）。 \\
Higgs 丛 / Higgs 场 & Higgs bundle / Higgs field & 全纯丛 $E$ 连同 $\phi\in H^{0}(C,K\otimes\ad(E))$；$\phi$ 称 Higgs 场。 \\
Hitchin 方程 & Hitchin's equations & $F-\phi\wedge\phi=0$、$D\phi=D\star\phi=0$；解模空间 $\MH$ 超凯勒。 \\
Hitchin 纤维化 & Hitchin fibration & $\MH\to B$，$(E,\phi)\mapsto\phi$ 的特征多项式；纤维是阿贝尔簇。 \\
谱曲线 & spectral curve & 特征方程 $\det(\lambda-\phi)=0$ 在 $K$ 的全纯线丛中所定义的曲线，用于把 Hitchin 纤维实现为 Prym 簇。 \\
可积系统 & integrable system & 此处指复意义下的 Liouville 可积：$\MH\to B$ 的纤维是环面、$B$ 是作用量变量空间。 \\
特殊拉格朗日子流形 & special Lagrangian submanifold & $\omega|_L=0$ 且 $\operatorname{Im}(e^{i\vartheta}\Omega)|_L=0$ 的实子流形；在 Kähler 流形中面积最小。 \\
T-对偶 & T-duality & 环面纤维化上的对偶：把每个纤维换成对偶环面，底不变。 \\
SYZ 猜想 & SYZ conjecture & 镜像对称 = 同一底上互为对偶的特殊拉格朗日环面纤维化。 \\
广义复几何 & generalized complex geometry & 在 $TM\oplus T^{*}M$ 上定义复结构，把复结构与辛结构放在同一个 $S^{2}$ 族中插值。 \\
D-模 & D-module & 微分算子层 $\mathcal{D}$ 上的（准）凝聚模；是带平坦联络的丛的无穷小版本。 \\
Hecke 算子 & Hecke operator & 在一点处对丛作基本修改所得的对应，作用在 D-模上。 \\
Hecke 特征层 & Hecke eigensheaf & 在 Hecke 算子下按张量倍（各不可约表示）作用的 D-模。 \\
膜 & brane & 拓扑场论中的边界条件；A-膜依赖辛结构，B-膜依赖复结构。 \\
稳定 / 半稳定 & stable / semistable & 按斜率（$\deg/\operatorname{rk}$）定义的丛的稳定性；稳定 $\Leftrightarrow$ 存在 Hitchin 解。 \\
分歧 & ramification & 平坦丛在穿刺点处有非平凡单值性的情形。 \\
驯顺分歧 & tame ramification & 单极点（留数可能非半单）的情形。 \\
野分歧 & wild ramification & 极点阶数大于 1 的情形；需用扩张单值性刻画。 \\
穿刺点 & puncture & 从 $C$ 中挖去的点 $p$；$C'=C\setminus p$ 上允许有单值性。 \\
单值性 & monodromy & 绕闭链平行移动所得的变换；对平坦丛即基本群表示。 \\
扩张单值性 & extended monodromy & 普通单值性连同 Stokes 矩阵，是非正则奇点的完整拓扑数据。 \\
Stokes 矩阵 & Stokes matrix & 跨过 Stokes 射线时渐近展开所乘的常数矩阵。 \\
等单值形变 & isomonodromic deformation & 保持（扩张）单值数据不变、只改变高阶极点系数的形变。 \\
辫群 & braid group & $\mathfrak{t}_{\C}^{\mathrm{reg}}$（或其 Weyl 商）的基本群；控制 $T_{n}$ 的绕行单值性。 \\
映射类群 & mapping class group & $C$ 的定向保持微分同胚模同伦；与辫群一起作用在野分歧的模空间上。 \\
叠 & stack & 局部可描述为商的空间；承载不稳定对象，使 Hecke 算子可作用。 \\
模叠 & moduli stack & 参数化对象的叠（如 $\Bun_{G}(C)$），与粗模空间相对。 \\
辛商 / 超凯勒商 & symplectic quotient / hyper-Kähler quotient & 取矩映射零点再模去群作用；后者同时取三个矩映射。 \\
矩映射 & moment map & Hamilton 群作用的生成元；其平方项出现在 $\sigma$ 模型作用量中。 \\
Bogomolny 方程 & Bogomolny equations & 三维规范场论中的单极子方程 $F=\star d_{A}\phi$。 \\
Wilson / 't Hooft 线算子 & Wilson / 't Hooft line operators & 沿一条曲线的电 / 磁探针；在对偶下互换。 \\
Montonen–Olive 对偶 & Montonen–Olive duality & $\mathcal{N}=4$ 理论的 S-对偶：$g\leftrightarrow1/g$、$G\leftrightarrow{}^{L}G$。 \\
Langlands 对偶群 & Langlands dual group & 根数据取对偶所得的群 ${}^{L}G$；环面情形即对偶环面。 \\
$\sigma$ 模型 & sigma model & 从世界面 $\Sigma$ 到目标 $X$ 的映射的量子场论。 \\
\end{longtable}

\subsection*{表 A2　记号}

\begin{longtable}{p{0.34\textwidth}p{0.60\textwidth}}
\toprule 记号 & 含义 \\ \midrule \endfirsthead
\toprule 记号 & 含义 \\ \midrule \endhead
\bottomrule \endfoot
$C$ & 紧致无边界定向二维流形（黎曼面），亏格 $g\ge2$ \\
$\YC(G,C)$ & 平坦 $G_{\C}$-丛模空间（表示簇） \\
$\MH(G,C)$ & Hitchin 方程解的模空间，超凯勒 \\
$B$ & Hitchin 底，$\bigoplus_{i}H^{0}(C,K^{d_{i}})$ \\
$\Omega$ & 复辛形式；$\Omega_{I}=\omega_{J}+i\omega_{K}$ \\
$\omega=\operatorname{Im}\Omega$ & A-模型所用的实辛形式 \\
$I,J,K$ & 超凯勒三复结构，$IJ=K$ \\
$\star$ & 复结构决定的 Hodge 星算子 \\
$\mathcal{A}=A+i\phi$ & 复联络（原文排作粗体 $\bm{A}$ 的地方） \\
$\mathfrak{t},\ \mathfrak{t}_{\C}^{\mathrm{reg}}$ & 极大环面的李代数 / 正则元开集 \\
$\alpha,\beta,\gamma$ & Simpson 参数（取值于 $\mathfrak{t}$），刻画穿刺点附近的奇性 \\
$U=\exp(-2\pi(\alpha-i\gamma))$ & 绕穿刺点的 holonomy（符号约定见译注 4.5） \\
$\eta$ & 量子参数（$\theta$-角的伙伴），A-模型侧的复化 Kähler 参数 \\
$T_{1},\dots,T_{n}$ & 野分歧极点系数，$T_{n}\in\mathfrak{t}_{\C}^{\mathrm{reg}}\otimes K_{p}^{n-1}$ \\
$B(G)$ & $G$ 的辫群 \\
$\Bun_{G}(C)$ & 所有 $G_{\C}$-丛的模叠 \\
$M(G,C)$ / $M_{\mathrm{st}}(G,C)$ & 半稳定 / 严格稳定丛的模空间 \\
${}^{L}G,\ {}^{L}\mathcal{C},\ {}^{L}T_{1}$ & 对偶群及对偶侧的相应对象 \\
\end{longtable}
